%%%%%%%%%%%%%%%%
%%% deut 31 %%%
%%%%%%%%%%%%%%%%

\Chapter{31}

\Verse{1}
{
	\hDtrA{וַיֵּלֶךְ מֹשֶׁה וַיְדַבֵּר אֶת הַדְּבָרִים הָאֵלֶּה אֶל כּל יִשְׂרָאֵל׃}
}
{
	\eDtrA{
		AND HE WENT 31 And Moses went and spoke these things to
		all Israel.
	}
}
{
%	\footnote{
%		And Moses went. Went where? The Septuagint and Qumran texts have “And
%		Moses finished speaking all these things”—reading Hebrew
%		(finished) rather than  (went), reversing the last two letters.
%		This makes better sense. It also adds a wordplay on his finishing
%		() in v. 1 and his saying “I’m not able () to go out and
%		come in anymore” in v. 2.
%	}
}

\Verse{2}
{
	\hDtrA{וַיֹּאמֶר אֲלֵהֶם בֶּן מֵאָה וְעֶשְׂרִים שָׁנָה אָנֹכִי הַיּוֹם לֹא אוּכַל עוֹד לָצֵאת וְלָבוֹא וַיהֹוָה אָמַר אֵלַי לֹא תַעֲבֹר אֶת הַיַּרְדֵּן הַזֶּה׃}
}
{
	\eDtrA{
		And he said to them, “I’m a hundred twenty years old
		today. I’m not able to go out and come in anymore. And
		YHWH said to me, ‘You shall not cross this Jordan.’
	}
}
{
}

\Verse{3}
{
	\hDtrA{יְהֹוָה אֱלֹהֶיךָ הוּא עֹבֵר לְפָנֶיךָ הוּא יַשְׁמִיד אֶת הַגּוֹיִם הָאֵלֶּה מִלְּפָנֶיךָ וִירִשְׁתָּם יְהוֹשֻׁעַ הוּא עֹבֵר לְפָנֶיךָ כַּאֲשֶׁר דִּבֶּר יְהֹוָה׃}
}
{
	\eDtrA{
		YHWH, your {\God}: He is crossing in front of you. He’ll
		destroy these nations in front of you, and you’ll
		dispossess them. Joshua: he is crossing in front of you,
		as YHWH has spoken.
	}
}
{
}

\Verse{4}
{
	\hDtrA{וְעָשָׂה יְהֹוָה לָהֶם כַּאֲשֶׁר עָשָׂה לְסִיחוֹן וּלְעוֹג מַלְכֵי הָאֱמֹרִי וּלְאַרְצָם אֲשֶׁר הִשְׁמִיד אֹתָם׃}
}
{
	\eDtrA{
		And YHWH will do to them as He did to Sihon and to Og, the
		kings of the Amorites, and to their land, that He
		destroyed them.
	}
}
{
}

\Verse{5}
{
	\hDtrA{וּנְתָנָם יְהֹוָה לִפְנֵיכֶם וַעֲשִׂיתֶם לָהֶם כְּכל הַמִּצְוָה אֲשֶׁר צִוִּיתִי אֶתְכֶם׃}
}
{
	\eDtrA{
		And YHWH will put them in front of you, and you shall do
		to them according to all of the commandment that I’ve
		commanded you.
	}
}
{
}

\Verse{6}
{
	\hDtrA{חִזְקוּ וְאִמְצוּ אַל תִּירְאוּ וְאַל תַּעַרְצוּ מִפְּנֵיהֶם כִּי יְהֹוָה אֱלֹהֶיךָ הוּא הַהֹלֵךְ עִמָּךְ לֹא יַרְפְּךָ וְלֹא יַעַזְבֶךָּ׃}
}
{
	\eDtrA{
		Be strong and be bold. Don’t be afraid and don’t be scared
		in front of them, because YHWH, your {\God}: He is the one
		going with you. He won’t let you down and won’t leave
		you.”
	}
}
{
}

\Verse{7}
{
	\hDtrA{וַיִּקְרָא מֹשֶׁה לִיהוֹשֻׁעַ וַיֹּאמֶר אֵלָיו לְעֵינֵי כל יִשְׂרָאֵל חֲזַק וֶאֱמָץ כִּי אַתָּה תָּבוֹא אֶת הָעָם הַזֶּה אֶל הָאָרֶץ אֲשֶׁר נִשְׁבַּע יְהֹוָה לַאֲבֹתָם לָתֵת לָהֶם וְאַתָּה תַּנְחִילֶנָּה אוֹתָם׃}
}
{
	\eDtrA{
		And Moses called Joshua and said to him before the eyes of
		all Israel, “Be strong and be bold, because you will come
		with this people to the land that YHWH swore to their
		fathers to give to them, and you will get it for them as a
		legacy.
	}
}
{
%	\footnote{
%		Be strong and be bold … you will get it for them as a legacy. This,
%		finally, is the culmination of what Moses told the people in chapter 3.
%		He had told them that God was angry at him, that he was not permitted to
%		enter the land, and that God told him to “command Joshua and strengthen
%		him and make him bold, because … he will get them the land that you’ll
%		see as a legacy” (3:28). Moses has now instructed the people in the laws
%		and exhorted them to be faithful to the covenant, and so he turns to the
%		task that he was assigned: he passes the leadership to Joshua. Note that
%		it is not just that Moses dies and Joshua replaces him. Rather, during
%		his lifetime, publicly, Moses himself identifies Joshua as his
%		successor, charges him with his task, and encourages him to be
%		successful. Although one can easily imagine Moses envying Joshua, who is
%		replacing him for the fulfillment of the promises, Moses acts in
%		Joshua’s favor: preparing him and enhancing his stature in the people’s
%		eyes. One who begrudges his successor does not understand time and
%		history.
%	}
}

\Verse{8}
{
	\hDtrA{וַיהֹוָה הוּא הַהֹלֵךְ לְפָנֶיךָ הוּא יִהְיֶה עִמָּךְ לֹא יַרְפְּךָ וְלֹא יַעַזְבֶךָּ לֹא תִירָא וְלֹא תֵחָת׃}
}
{
	\eDtrA{
		And YHWH: He is the one who is going in front of you. He
		will be with you. He won’t let you down and won’t leave
		you. You shall not fear, and you shall not be dismayed.”
	}
}
{
}

\Verse{9}
{
	\hDtrA{וַיִּכְתֹּב מֹשֶׁה אֶת הַתּוֹרָה הַזֹּאת וַיִּתְּנָהּ אֶל הַכֹּהֲנִים בְּנֵי לֵוִי הַנֹּשְׂאִים אֶת אֲרוֹן בְּרִית יְהֹוָה וְאֶל כּל זִקְנֵי יִשְׂרָאֵל׃}
}
{
	\eDtrA{
		And Moses wrote this instruction and gave it to the
		priests, sons of Levi, who were carrying the ark of YHWH’s
		covenant and to all of Israel’s elders.
	}
}
{
%	\footnote{
%		And Moses wrote this instruction (and 31:24. “And it was when Moses
%		finished writing the words of this instruction on a scroll”). This
%		“instruction” (Hebrew tôrh) does not refer to the entire five
%		books, which came to be known as the Torah long after this chapter was
%		written. Those who began claiming that Moses was the author of the
%		entire Torah created much confusion and misinformation by claiming more
%		for Moses than the Torah itself ascribed to him. It became a firmly
%		established doctrine by the rabbinic period, and many great teachers,
%		including Rashi and Ramban, followed it. Ibn Ezra raised subtle doubts
%		about it, Spinoza openly challenged it, and Jewish, as well as
%		Christian, scholars came to reject it in the nineteenth and twentieth
%		centuries. Orthodox Jews still accept it. The instruction that Moses
%		writes here appears to be most of the law code that begins in
%		Deuteronomy 12 and ends at 26:15. It may also be understood to include
%		the list of blessings and curses in Deuteronomy 28. That concluding list
%		of curses especially would account for the strong reaction that is
%		ascribed to King Josiah when he hears the words of the scroll of the
%		Torah centuries later (see the comment on 31:26).
%	}
}

\Verse{10}
{
	\hDtrA{וַיְצַו מֹשֶׁה אוֹתָם לֵאמֹר מִקֵּץ שֶׁבַע שָׁנִים בְּמֹעֵד שְׁנַת הַשְּׁמִטָּה בְּחַג הַסֻּכּוֹת׃}
}
{
	\eDtrA{
		And Moses commanded them, saying, “At the end of seven
		years, at the appointed time of the year of the remission,
		on the Festival of Booths,
	}
}
{
%	\footnote{
%		seven years. The biblical requirement was to read “this torah” (the law
%		code of Deuteronomy) once every seven years, on the holiday of Sukkot.
%		Later the Jewish practice became to read the entire Torah (all of the
%		Five Books of Moses), once every year or every three years, and in
%		weekly portions rather than all at once. Footnote 31:10. the remission.
%		See Deut 15:1ff.
%	}
}

\Verse{11}
{
	\hDtrA{בְּבוֹא כל יִשְׂרָאֵל לֵרָאוֹת אֶת פְּנֵי יְהֹוָה אֱלֹהֶיךָ בַּמָּקוֹם אֲשֶׁר יִבְחָר תִּקְרָא אֶת הַתּוֹרָה הַזֹּאת נֶגֶד כּל יִשְׂרָאֵל בְּאזְנֵיהֶם׃}
}
{
	\eDtrA{
		when all Israel comes to appear before YHWH, your {\God}, in
		the place that He will choose, you shall read this
		instruction in front of all Israel in their ears.
	}
}
{
}

\Verse{12}
{
	\hDtrA{הַקְהֵל אֶת הָעָם הָאֲנָשִׁים וְהַנָּשִׁים וְהַטַּף וְגֵרְךָ אֲשֶׁר בִּשְׁעָרֶיךָ לְמַעַן יִשְׁמְעוּ וּלְמַעַן יִלְמְדוּ וְיָרְאוּ אֶת יְהֹוָה אֱלֹהֵיכֶם וְשָׁמְרוּ לַעֲשׂוֹת אֶת כּל דִּבְרֵי הַתּוֹרָה הַזֹּאת׃}
}
{
	\eDtrA{
		Assemble the people—the men and the women and the infants
		and your alien who is in your gates—so they will listen
		and so they will learn and will fear YHWH, your {\God}, and
		they will be watchful to do all the words of this
		instruction.
	}
}
{
%	\footnote{
%		your alien. Five times in the Torah, laws about aliens in Israel say
%		“your alien” (compare “your poor” in the Covenant Code; Exod
%		23:6)—including in the Decalogue in both Exodus and Deuteronomy.
%	}
}

\Verse{13}
{
	\hDtrA{וּבְנֵיהֶם אֲשֶׁר לֹא יָדְעוּ יִשְׁמְעוּ וְלָמְדוּ לְיִרְאָה אֶת יְהֹוָה אֱלֹהֵיכֶם כּל הַיָּמִים אֲשֶׁר אַתֶּם חַיִּים עַל הָאֲדָמָה אֲשֶׁר אַתֶּם עֹבְרִים אֶת הַיַּרְדֵּן שָׁמָּה לְרִשְׁתָּהּ׃}
}
{
	\eDtrA{
		And their children who have not known will listen and
		learn to fear YHWH, your {\God}, all the days that you’re
		living on the land to which you’re crossing the Jordan to
		take possession of it.”
	}
}
{
}

\Verse{14}
{
	\hE{וַיֹּאמֶר יְהֹוָה אֶל מֹשֶׁה הֵן קָרְבוּ יָמֶיךָ לָמוּת קְרָא אֶת יְהוֹשֻׁעַ וְהִתְיַצְּבוּ בְּאֹהֶל מוֹעֵד וַאֲצַוֶּנּוּ וַיֵּלֶךְ מֹשֶׁה וִיהוֹשֻׁעַ וַיִּתְיַצְּבוּ בְּאֹהֶל מוֹעֵד׃}
}
{
	\eE{
		And YHWH said to Moses, “Here, your days to die have come
		close. Call Joshua, and stand up in the Tent of Meeting,
		and I’ll command him.” And Moses and Joshua went and stood
		up in the Tent of Meeting.
	}
}
{
}

\Verse{15}
{
	\hE{וַיֵּרָא יְהֹוָה בָּאֹהֶל בְּעַמּוּד עָנָן וַיַּעֲמֹד עַמּוּד הֶעָנָן עַל פֶּתַח הָאֹהֶל׃}
}
{
	\eE{
		And YHWH appeared in the tent in a column of cloud, and
		the column of cloud stood at the entrance of the tent.
	}
}
{
}

\Verse{16}
{
	\hDtrB{וַיֹּאמֶר יְהֹוָה אֶל מֹשֶׁה הִנְּךָ שֹׁכֵב עִם אֲבֹתֶיךָ וְקָם הָעָם הַזֶּה וְזָנָה אַחֲרֵי אֱלֹהֵי נֵכַר הָאָרֶץ אֲשֶׁר הוּא בָא שָׁמָּה בְּקִרְבּוֹ וַעֲזָבַנִי וְהֵפֵר אֶת בְּרִיתִי אֲשֶׁר כָּרַתִּי אִתּוֹ׃}
}
{
	\eDtrB{
		And YHWH said to Moses, “Here, when you’re lying with your
		fathers, this people will get up and whore after foreign
		gods of the land into which it is coming, and it will
		leave me and break my covenant that I’ve made with it.
	}
}
{
%	\footnote{
%		it. God repeatedly speaks of the people in the singular in this
%		passage—and quotes the people as speaking of itself in the singular as
%		well: “not present in me … these evils have found me.” The hiding of the
%		face of God, which is predicted here, is understood to be the fate of
%		the entire people, not of individuals. Divine hiddenness is not an
%		individual’s unique experience, and it is not based on an individual’s
%		behavior. It is the experience of a community and possibly of an era. At
%		the time that I am writing this, it is such an era.
%	}
}

\Verse{17}
{
	\hDtrB{וְחָרָה אַפִּי בוֹ בַיּוֹם הַהוּא וַעֲזַבְתִּים וְהִסְתַּרְתִּי פָנַי מֵהֶם וְהָיָה לֶאֱכֹל וּמְצָאֻהוּ רָעוֹת רַבּוֹת וְצָרוֹת וְאָמַר בַּיּוֹם הַהוּא הֲלֹא עַל כִּי אֵין אֱלֹהַי בְּקִרְבִּי מְצָאוּנִי הָרָעוֹת הָאֵלֶּה׃}
}
{
	\eDtrB{
		And my anger will flare at it on that day, and I’ll leave
		them, and I’ll hide my face from them, and it will become
		prey, and many bad things and troubles will find it. And
		it will say on that day, ‘Isn’t it because my {\God} is not
		present in me that these evils have found me?’
	}
}
{
%	\footnote{
%		I’ll hide my face. These words come twice in the passage that contains
%		the last words that God speaks to Moses before summoning him to die.
%		They occur in the song that God tells Moses to teach the people as a
%		witness for the future as well (32:20), and this expression occurs
%		nearly thirty times more in the Tanak after that. These words here
%		predict a time in which God will be hidden. It is more frightening than
%		divine punishment: It is one thing for a parent to punish a child; it is
%		far worse if the parent becomes hidden from the child, so the child does
%		not know if the parent is present or not. The child feels vulnerable,
%		unprotected. The child cries for the parent, but there is no answer. It
%		is terrifying.
%	}
}

\Verse{18}
{
	\hDtrB{וְאָנֹכִי הַסְתֵּר אַסְתִּיר פָּנַי בַּיּוֹם הַהוּא עַל כּל הָרָעָה אֲשֶׁר עָשָׂה כִּי פָנָה אֶל אֱלֹהִים אֲחֵרִים׃}
}
{
	\eDtrB{
		But I: I’ll hide my face on that day over all the bad that
		it has done, because it turned to other gods.
	}
}
{
%	\footnote{
%		it. God repeatedly speaks of the people in the singular in this
%		passage—and quotes the people as speaking of itself in the singular as
%		well: “not present in me … these evils have found me.” The hiding of the
%		face of God, which is predicted here, is understood to be the fate of
%		the entire people, not of individuals. Divine hiddenness is not an
%		individual’s unique experience, and it is not based on an individual’s
%		behavior. It is the experience of a community and possibly of an era. At
%		the time that I am writing this, it is such an era. Footnote 31:18. I’ll
%		hide my face. These words come twice in the passage that contains the
%		last words that God speaks to Moses before summoning him to die. They
%		occur in the song that God tells Moses to teach the people as a witness
%		for the future as well (32:20), and this expression occurs nearly thirty
%		times more in the Tanak after that. These words here predict a time in
%		which God will be hidden. It is more frightening than divine punishment:
%		It is one thing for a parent to punish a child; it is far worse if the
%		parent becomes hidden from the child, so the child does not know if the
%		parent is present or not. The child feels vulnerable, unprotected. The
%		child cries for the parent, but there is no answer. It is terrifying.
%		The expression hestr pnîm, the hiding of the face, became
%		a known idiom in rabbinic literature, and in theological thinking to
%		this day. I have written about how it develops through the Bible (in The
%		Hidden Face of God). The Tanak’s story moves from a time of
%		extraordinary divine involvement—creation, flood, personal contact,
%		splitting sea, hearing God’s voice at Sinai—to a time in which humans
%		are left on their own. God is not pictured as speaking or appearing in
%		the books of Ezra and Nehemiah and is never mentioned in Esther. The
%		prediction at the end of the Torah comes true at the end of the Tanak.
%		Humans are left not knowing if God exists or not. Thus the Tanak tells a
%		story of a development from a world of direct communication with God to
%		the world we have known ever since: in which the existence of God is a
%		matter of faith, or doubt, or search.
%	}
}

\Verse{19}
{
	\hDtrB{וְעַתָּה כִּתְבוּ לָכֶם אֶת הַשִּׁירָה הַזֹּאת וְלַמְּדָהּ אֶת בְּנֵי יִשְׂרָאֵל שִׂימָהּ בְּפִיהֶם לְמַעַן תִּהְיֶה לִּי הַשִּׁירָה הַזֹּאת לְעֵד בִּבְנֵי יִשְׂרָאֵל׃}
}
{
	\eDtrB{
		So now write this song and teach it to the children of
		Israel. Set it in their mouths, so this song will become a
		witness for me among the children of Israel.
	}
}
{
}

\Verse{20}
{
	\hDtrB{כִּי אֲבִיאֶנּוּ אֶל הָאֲדָמָה אֲשֶׁר נִשְׁבַּעְתִּי לַאֲבֹתָיו זָבַת חָלָב וּדְבַשׁ וְאָכַל וְשָׂבַע וְדָשֵׁן וּפָנָה אֶל אֱלֹהִים אֲחֵרִים וַעֲבָדוּם וְנִאֲצוּנִי וְהֵפֵר אֶת בְּרִיתִי׃}
}
{
	\eDtrB{
		When I’ll bring it to the land that I swore to its
		fathers, flowing with milk and honey, and it will eat and
		be full and get fat and turn to other gods, and they will
		serve them and reject me, and it will break my covenant,
	}
}
{
}

\Verse{21}
{
	\hDtrB{וְהָיָה כִּי תִמְצֶאןָ אֹתוֹ רָעוֹת רַבּוֹת וְצָרוֹת וְעָנְתָה הַשִּׁירָה הַזֹּאת לְפָנָיו לְעֵד כִּי לֹא תִשָּׁכַח מִפִּי זַרְעוֹ כִּי יָדַעְתִּי אֶת יִצְרוֹ אֲשֶׁר הוּא עֹשֶׂה הַיּוֹם בְּטֶרֶם אֲבִיאֶנּוּ אֶל הָאָרֶץ אֲשֶׁר נִשְׁבַּעְתִּי׃}
}
{
	\eDtrB{
		then it will be, when many bad things and troubles will
		find it, that this song will testify as a witness in front
		of it, because it won’t be forgotten from its seed’s
		mouth. Because I know its inclination that it is doing
		today even before I bring it to the land that I swore.”
	}
}
{
%	\footnote{
%		I know its inclination. Again an essential point from the Torah’s
%		beginning returns at its conclusion. The flood story starts and ends
%		with God’s recognition that humans’ inclination is toward doing bad: “I
%		won’t curse the ground on account of humankind again, for the
%		inclination of the human heart is bad from their youth” (Gen 6:5; 8:21).
%		The word “inclination” (yer) then does not occur again in
%		the Torah until here. God states that He knows, even before the people
%		enter the land, that this will be a problem, and so He has them learn a
%		song that refers to their rebellion and to the hiding of the face of God
%		so that they will know that they were warned.
%	}
}

\Verse{22}
{
	\hDtrB{וַיִּכְתֹּב מֹשֶׁה אֶת הַשִּׁירָה הַזֹּאת בַּיּוֹם הַהוּא וַיְלַמְּדָהּ אֶת בְּנֵי יִשְׂרָאֵל׃}
}
{
	\eDtrB{
		And Moses wrote this song on that day, and he taught it to
		the children of Israel.
	}
}
{
}

\Verse{23}
{
	\hE{וַיְצַו אֶת יְהוֹשֻׁעַ בִּן נוּן וַיֹּאמֶר חֲזַק וֶאֱמָץ כִּי אַתָּה תָּבִיא אֶת בְּנֵי יִשְׂרָאֵל אֶל הָאָרֶץ אֲשֶׁר נִשְׁבַּעְתִּי לָהֶם וְאָנֹכִי אֶהְיֶה עִמָּךְ׃}
}
{
	\eE{
		And He commanded Joshua, son of Nun, and said, “Be strong
		and bold, because you will bring the children of Israel to
		the land that I swore to them, and I shall be with you.”
	}
}
{
}

\Verse{24}
{
	\hDtrA{וַיְהִי כְּכַלּוֹת מֹשֶׁה לִכְתֹּב אֶת דִּבְרֵי הַתּוֹרָה הַזֹּאת עַל סֵפֶר עַד תֻּמָּם׃}
}
{
	\eDtrA{
		And it was when Moses finished writing the words of this
		instruction on a scroll to their end,
	}
}
{
}

\Verse{25}
{
	\hDtrA{וַיְצַו מֹשֶׁה אֶת הַלְוִיִּם נֹשְׂאֵי אֲרוֹן בְּרִית יְהֹוָה לֵאמֹר׃}
}
{
	\eDtrA{
		and Moses commanded the Levites, who carried the ark of
		the covenant of YHWH, saying,
	}
}
{
}

\Verse{26}
{
	\hDtrA{לָקֹחַ אֵת סֵפֶר הַתּוֹרָה הַזֶּה וְשַׂמְתֶּם אֹתוֹ מִצַּד אֲרוֹן בְּרִית יְהֹוָה אֱלֹהֵיכֶם וְהָיָה שָׁם בְּךָ לְעֵד׃}
}
{
	\eDtrA{
		“Take this scroll of instruction and set it at the side of
		the ark of the covenant of YHWH, your {\God}, and it will be
		there for you as a witness.
	}
}
{
%	\footnote{
%		it will be there for you as a witness. The scroll that Moses has written
%		stays there beside the ark for six hundred years until the reign of King
%		Josiah. Apparently it has long been ignored by then. The priest Hilkiah
%		says there, “I’ve found the scroll of instruction (sper
%		hattrh) in the house of YHWH.” When it is read to King
%		Josiah, he tears his clothes and says, “YHWH’s fury is great that has
%		ignited at us because our fathers didn’t listen to this scroll’s words,
%		to act according to everything that’s written about us!” (2 Kings
%		22:8–13). Especially in light of the curses in Deuteronomy 28, Josiah’s
%		reaction is understandable. And so it is understood that these laws go
%		ignored for generations and then are rediscovered. Josiah has the scroll
%		read aloud to the entire people, and he institutes a renewal of the
%		covenant and a tremendous religious reform. The narrative concludes:
%		“And before him there was not a king like him, who came back to YHWH
%		with all his heart and with all his soul and with all his might,
%		according to all of Moses’ instruction, and after him none rose like
%		him” (2 Kings 23:25). All the kings in the books of 1 and 2 Kings are
%		measured by whether they fulfilled the laws of Deuteronomy, and Josiah
%		is rated the highest. Above all, Josiah institutes the practice of
%		public reading of the Torah, which prevails to the present day. And, to
%		this day, the measure of a leader among the Jews must be the degree to
%		which he or she lives by the Torah and draws others to do so as well.
%	}
}

\Verse{27}
{
	\hDtrA{כִּי אָנֹכִי יָדַעְתִּי אֶת מֶרְיְךָ וְאֶת ערְפְּךָ הַקָּשֶׁה הֵן בְּעוֹדֶנִּי חַי עִמָּכֶם הַיּוֹם מַמְרִים הֱיִתֶם עִם יְהֹוָה וְאַף כִּי אַחֲרֵי מוֹתִי׃}
}
{
	\eDtrA{
		Because I know your rebellion and your hard neck. Here,
		while I’m still alive with you today, you’ve been
		rebelling at YHWH, so how much more after my death!
	}
}
{
}

\Verse{28}
{
	\hDtrB{הַקְהִילוּ אֵלַי אֶת כּל זִקְנֵי שִׁבְטֵיכֶם וְשֹׁטְרֵיכֶם וַאֲדַבְּרָה בְאזְנֵיהֶם אֵת הַדְּבָרִים הָאֵלֶּה וְאָעִידָה בָּם אֶת הַשָּׁמַיִם וְאֶת הָאָרֶץ׃}
}
{
	\eDtrB{
		Assemble all the elders of your tribes and your officers
		to me so I may speak these things in their ears and call
		the skies and the earth to witness regarding them,
	}
}
{
%	\footnote{
%		call the skies and the earth to witness. Moses will do this in the
%		opening words of the song (32:1): “Listen, skies, so I may speak; and
%		let the earth hear what my mouth says.” Ancient Near Eastern treaties
%		include a summoning of witnesses, usually the gods of the parties to the
%		agreement. The biblical covenant cannot include pagan gods as witnesses,
%		but it must have eternal witnesses since the covenant is meant to last
%		forever. Therefore, the earth and sky are called on as witnesses. This
%		summoning of earth and sky rather than gods in an early poem is further
%		evidence that monotheism existed early in Israelite history.
%	}
}

\Verse{29}
{
	\hDtrB{כִּי יָדַעְתִּי אַחֲרֵי מוֹתִי כִּי הַשְׁחֵת תַּשְׁחִתוּן וְסַרְתֶּם מִן הַדֶּרֶךְ אֲשֶׁר צִוִּיתִי אֶתְכֶם וְקָרָאת אֶתְכֶם הָרָעָה בְּאַחֲרִית הַיָּמִים כִּי תַעֲשׂוּ אֶת הָרַע בְּעֵינֵי יְהֹוָה לְהַכְעִיסוֹ בְּמַעֲשֵׂה יְדֵיכֶם׃}
}
{
	\eDtrB{
		because I know, after my death, that you’ll be corrupted,
		and you’ll turn from the way that I’ve commanded you. And
		the bad thing will happen to you in the future days when
		you’ll do what is bad in YHWH’s eyes to provoke Him with
		your hands’ work.”
	}
}
{
%	\footnote{
%		you’ll be corrupted. Another term from the Torah’s beginning returns at
%		its end: the Hiphil of the root t with the sense of humans
%		becoming corrupt was used to introduce the flood: “All flesh had
%		corrupted its way on the earth” (Gen 6:12). It does not occur again with
%		this meaning until Deut 4:16 and then again here, where Moses says he
%		knows that the people will someday become corrupted. That is, he
%		understands that the susceptibility to corruption has been part of human
%		nature from the beginning.
%	}
}

\Verse{30}
{
	\hDtrB{וַיְדַבֵּר מֹשֶׁה בְּאזְנֵי כּל קְהַל יִשְׂרָאֵל אֶת דִּבְרֵי הַשִּׁירָה הַזֹּאת עַד תֻּמָּם׃}
}
{
	\eDtrB{
		And Moses spoke in the ears of all the community of Israel
		the words of this song to their end:
	}
}
{
}
